\documentclass[conference]{IEEEtran}
\usepackage{amsmath,amssymb,amsthm}
\usepackage{graphicx}
\usepackage{booktabs}
\usepackage{algorithm}
\usepackage{algpseudocode}
\usepackage{xcolor}
\usepackage{url}
\usepackage[hidelinks]{hyperref}
\usepackage{cite}

% AUTO-GENERATED by experiments/analyze.py -- do not edit
\newcommand{\abRoughOne}{22.3}
\newcommand{\abRoughZero}{17.9}
\newcommand{\abRoughZeropFive}{20.3}
\newcommand{\abRoughZeropSevenFive}{21.3}
\newcommand{\abRoughZeropTwoFive}{19.4}
\newcommand{\abSlopeOne}{7.9}
\newcommand{\abSlopeZero}{13.5}
\newcommand{\abSlopeZeropFive}{8.6}
\newcommand{\abSlopeZeropSevenFive}{8.1}
\newcommand{\abSlopeZeropTwoFive}{8.8}
\newcommand{\climbRelA}{0.0}
\newcommand{\climbRelAHC}{-9.3}
\newcommand{\climbRelTA}{-37.6}
\newcommand{\climbRelTAS}{-40.2}
\newcommand{\climbRelTheta}{-4.7}
\newcommand{\costRelA}{0.0}
\newcommand{\costRelAHC}{-4.0}
\newcommand{\costRelTA}{-50.2}
\newcommand{\costRelTAS}{-48.2}
\newcommand{\costRelTAScraters}{-62.1}
\newcommand{\costRelTAShills}{-35.4}
\newcommand{\costRelTASridge}{-51.3}
\newcommand{\costRelTASrocky}{-45.3}
\newcommand{\costRelTheta}{-0.0}
\newcommand{\dTASAHCclimb}{-0.95}
\newcommand{\dTASAHCenergykj}{-0.55}
\newcommand{\dTASAHClengthxy}{+4.31}
\newcommand{\dTASAHCmeancost}{-0.125}
\newcommand{\dTASAHCpNineFiverough}{-4.5}
\newcommand{\dTASAHCpNineFiveslope}{-9.27}
\newcommand{\dTASAclimb}{-1.69}
\newcommand{\dTASAenergykj}{-1.24}
\newcommand{\dTASAlengthxy}{+5.38}
\newcommand{\dTASAmeancost}{-0.152}
\newcommand{\dTASApNineFiverough}{-7.1}
\newcommand{\dTASApNineFiveslope}{-12.30}
\newcommand{\dTASTAclimb}{-0.06}
\newcommand{\dTASTAenergykj}{-1.06}
\newcommand{\dTASTAlengthxy}{-4.41}
\newcommand{\dTASTAmeancost}{+0.010}
\newcommand{\dTASTApNineFiverough}{+0.4}
\newcommand{\dTASTApNineFiveslope}{+0.42}
\newcommand{\dTASThetaclimb}{-1.66}
\newcommand{\dTASThetaenergykj}{+0.22}
\newcommand{\dTASThetalengthxy}{+9.95}
\newcommand{\dTASThetameancost}{-0.155}
\newcommand{\dTASThetapNineFiverough}{-5.7}
\newcommand{\dTASThetapNineFiveslope}{-11.97}
\newcommand{\energyRelA}{0.0}
\newcommand{\energyRelAHC}{-1.4}
\newcommand{\energyRelTA}{-0.7}
\newcommand{\energyRelTAS}{-5.0}
\newcommand{\energyRelTAScraters}{-14.7}
\newcommand{\energyRelTAShills}{-0.7}
\newcommand{\energyRelTASridge}{-6.7}
\newcommand{\energyRelTASrocky}{-4.2}
\newcommand{\energyRelTheta}{-4.6}
\newcommand{\exposedA}{76}
\newcommand{\exposedAHC}{0}
\newcommand{\exposedAcraters}{83}
\newcommand{\exposedAhills}{63}
\newcommand{\exposedAridge}{97}
\newcommand{\exposedArocky}{60}
\newcommand{\exposedTA}{0}
\newcommand{\exposedTAS}{0}
\newcommand{\exposedTheta}{82}
\newcommand{\exposedThetacraters}{80}
\newcommand{\exposedThetahills}{77}
\newcommand{\exposedThetaridge}{90}
\newcommand{\exposedThetarocky}{80}
\newcommand{\lenRelA}{0.0}
\newcommand{\lenRelAHC}{0.1}
\newcommand{\lenRelTA}{11.8}
\newcommand{\lenRelTAS}{6.6}
\newcommand{\lenRelTAScraters}{7.2}
\newcommand{\lenRelTAShills}{6.3}
\newcommand{\lenRelTASridge}{9.6}
\newcommand{\lenRelTASrocky}{5.0}
\newcommand{\lenRelTheta}{-5.7}
\newcommand{\maxPitchA}{26.8}
\newcommand{\maxPitchAHC}{19.0}
\newcommand{\maxPitchTA}{11.7}
\newcommand{\maxPitchTAS}{11.4}
\newcommand{\maxPitchTheta}{25.3}
\newcommand{\maxSlopeA}{31.1}
\newcommand{\maxSlopeAHC}{23.7}
\newcommand{\maxSlopeTA}{12.4}
\newcommand{\maxSlopeTAS}{12.9}
\newcommand{\maxSlopeTheta}{33.2}
\newcommand{\nCPU}{4}
\newcommand{\nRedraw}{1}
\newcommand{\nSeeds}{10}
\newcommand{\nTest}{120}
\newcommand{\nTests}{24}
\newcommand{\nVal}{60}
\newcommand{\noiseCostAHCZeropOne}{0.267}
\newcommand{\noiseCostAHCZeropZeroFive}{0.287}
\newcommand{\noiseCostAHCZeropZeroTwo}{0.292}
\newcommand{\noiseCostTAZeropOne}{0.205}
\newcommand{\noiseCostTAZeropZeroFive}{0.180}
\newcommand{\noiseCostTAZeropZeroTwo}{0.154}
\newcommand{\noiseSteepAHCZeropOne}{0.10}
\newcommand{\noiseSteepAHCZeropZeroFive}{0.14}
\newcommand{\noiseSteepAHCZeropZeroTwo}{0.04}
\newcommand{\noiseSteepTAZeropOne}{0.03}
\newcommand{\noiseSteepTAZeropZeroFive}{0.00}
\newcommand{\noiseSteepTAZeropZeroTwo}{0.00}
\newcommand{\noiseSuccAHCZeropOne}{99.2}
\newcommand{\noiseSuccAHCZeropZeroFive}{100.0}
\newcommand{\noiseSuccAHCZeropZeroTwo}{100.0}
\newcommand{\noiseSuccTAZeropOne}{99.2}
\newcommand{\noiseSuccTAZeropZeroFive}{100.0}
\newcommand{\noiseSuccTAZeropZeroTwo}{100.0}
\newcommand{\pNineFiveSlopeA}{25.2}
\newcommand{\pNineFiveSlopeAHC}{19.7}
\newcommand{\pNineFiveSlopeTA}{8.1}
\newcommand{\pNineFiveSlopeTAS}{8.7}
\newcommand{\pNineFiveSlopeTheta}{24.9}
\newcommand{\pTASAHCclimb}{$<$0.001}
\newcommand{\pTASAHCenergykj}{$<$0.001}
\newcommand{\pTASAHClengthxy}{$<$0.001}
\newcommand{\pTASAHCmeancost}{$<$0.001}
\newcommand{\pTASAHCpNineFiverough}{$<$0.001}
\newcommand{\pTASAHCpNineFiveslope}{$<$0.001}
\newcommand{\pTASAclimb}{$<$0.001}
\newcommand{\pTASAenergykj}{$<$0.001}
\newcommand{\pTASAlengthxy}{$<$0.001}
\newcommand{\pTASAmeancost}{$<$0.001}
\newcommand{\pTASApNineFiverough}{$<$0.001}
\newcommand{\pTASApNineFiveslope}{$<$0.001}
\newcommand{\pTASTAclimb}{$<$0.001}
\newcommand{\pTASTAenergykj}{$<$0.001}
\newcommand{\pTASTAlengthxy}{$<$0.001}
\newcommand{\pTASTAmeancost}{$<$0.001}
\newcommand{\pTASTApNineFiverough}{$<$0.001}
\newcommand{\pTASTApNineFiveslope}{$<$0.001}
\newcommand{\pTASThetaclimb}{$<$0.001}
\newcommand{\pTASThetaenergykj}{0.973}
\newcommand{\pTASThetalengthxy}{$<$0.001}
\newcommand{\pTASThetameancost}{$<$0.001}
\newcommand{\pTASThetapNineFiverough}{$<$0.001}
\newcommand{\pTASThetapNineFiveslope}{$<$0.001}
\newcommand{\rTASAHCclimb}{-0.98}
\newcommand{\rTASAHCenergykj}{-0.46}
\newcommand{\rTASAHClengthxy}{+0.90}
\newcommand{\rTASAHCmeancost}{-1.00}
\newcommand{\rTASAHCpNineFiverough}{-0.98}
\newcommand{\rTASAHCpNineFiveslope}{-1.00}
\newcommand{\rTASAclimb}{-0.99}
\newcommand{\rTASAenergykj}{-0.67}
\newcommand{\rTASAlengthxy}{+0.93}
\newcommand{\rTASAmeancost}{-1.00}
\newcommand{\rTASApNineFiverough}{-0.99}
\newcommand{\rTASApNineFiveslope}{-1.00}
\newcommand{\rTASTAclimb}{-0.85}
\newcommand{\rTASTAenergykj}{-1.00}
\newcommand{\rTASTAlengthxy}{-1.00}
\newcommand{\rTASTAmeancost}{+1.00}
\newcommand{\rTASTApNineFiverough}{+0.58}
\newcommand{\rTASTApNineFiveslope}{+0.98}
\newcommand{\rTASThetaclimb}{-0.95}
\newcommand{\rTASThetaenergykj}{-0.00}
\newcommand{\rTASThetalengthxy}{+1.00}
\newcommand{\rTASThetameancost}{-1.00}
\newcommand{\rTASThetapNineFiverough}{-0.99}
\newcommand{\rTASThetapNineFiveslope}{-1.00}
\newcommand{\roughRockyA}{113}
\newcommand{\roughRockyAHC}{104}
\newcommand{\roughRockyTA}{43}
\newcommand{\roughRockyTAS}{43}
\newcommand{\roughRockyTheta}{116}
\newcommand{\steepMeanA}{8.5}
\newcommand{\steepMeanAHC}{0.0}
\newcommand{\steepMeanTA}{0.0}
\newcommand{\steepMeanTAS}{0.0}
\newcommand{\steepMeanTheta}{8.1}
\newcommand{\succA}{100.0}
\newcommand{\succAHC}{100.0}
\newcommand{\succTA}{100.0}
\newcommand{\succTAS}{100.0}
\newcommand{\succTheta}{100.0}
\newcommand{\sweepCostEight}{0.157}
\newcommand{\sweepCostFour}{0.175}
\newcommand{\sweepCostOne}{0.234}
\newcommand{\sweepCostOneSix}{0.150}
\newcommand{\sweepCostThreeTwo}{0.145}
\newcommand{\sweepCostTwo}{0.198}
\newcommand{\sweepCostZero}{0.295}
\newcommand{\sweepCostZeropFive}{0.260}
\newcommand{\sweepEnergyEight}{-6.8}
\newcommand{\sweepEnergyFour}{-8.6}
\newcommand{\sweepEnergyOne}{-9.3}
\newcommand{\sweepEnergyOneSix}{-5.0}
\newcommand{\sweepEnergyThreeTwo}{-3.0}
\newcommand{\sweepEnergyTwo}{-9.6}
\newcommand{\sweepEnergyZero}{-6.6}
\newcommand{\sweepEnergyZeropFive}{-8.7}
\newcommand{\sweepLenEight}{3.7}
\newcommand{\sweepLenFour}{1.5}
\newcommand{\sweepLenOne}{-2.4}
\newcommand{\sweepLenOneSix}{6.6}
\newcommand{\sweepLenThreeTwo}{8.5}
\newcommand{\sweepLenTwo}{-0.9}
\newcommand{\sweepLenZero}{-3.1}
\newcommand{\sweepLenZeropFive}{-3.1}
\newcommand{\timeA}{0.037}
\newcommand{\timeAFiveZero}{0.008}
\newcommand{\timeAHC}{0.035}
\newcommand{\timeAHCFiveZero}{0.008}
\newcommand{\timeAHCOneZeroZero}{0.036}
\newcommand{\timeAHCTwoZeroZero}{0.148}
\newcommand{\timeAOneZeroZero}{0.038}
\newcommand{\timeATwoZeroZero}{0.153}
\newcommand{\timeMaxA}{0.30}
\newcommand{\timeMaxAHC}{0.30}
\newcommand{\timeMaxTA}{1.49}
\newcommand{\timeMaxTAS}{2.17}
\newcommand{\timeMaxTheta}{0.79}
\newcommand{\timeTA}{0.180}
\newcommand{\timeTAFiveZero}{0.049}
\newcommand{\timeTAOneZeroZero}{0.180}
\newcommand{\timeTAS}{0.294}
\newcommand{\timeTASFiveZero}{0.087}
\newcommand{\timeTASOneZeroZero}{0.294}
\newcommand{\timeTASTwoZeroZero}{1.295}
\newcommand{\timeTATwoZeroZero}{0.823}
\newcommand{\timeTheta}{0.057}
\newcommand{\timeThetaFiveZero}{0.017}
\newcommand{\timeThetaOneZeroZero}{0.058}
\newcommand{\timeThetaTwoZeroZero}{0.170}
\newcommand{\valOverheadMain}{8.7}
\newcommand{\wMain}{16}
\newcommand{\wallMin}{7}

\newtheorem{proposition}{Proposition}
\newcommand{\TAS}{TA$^*$-S}
\newcommand{\TA}{TA$^*$}
\newcommand{\Astar}{A$^*$}
\newcommand{\HC}{A$^*$-HC}
\newcommand{\thmax}{\theta_{\max}}

\begin{document}

\title{Terrain-Aware A* with Robot-Grounded Traversability Costs:\\
A Reproducible Evaluation for Tracked UGVs on Uneven Terrain}

% TODO(author): fill in names and affiliations
\author{\IEEEauthorblockN{Author Name}
\IEEEauthorblockA{Affiliation\\ City, Country\\ email@example.com}}

\maketitle

\begin{abstract}
Shortest-path planners routinely send ground robots across slopes and rough
patches that a cost-aware planner would avoid at small extra distance. We
study a simple and fully specified terrain-aware planner for tracked unmanned
ground vehicles (UGVs): a traversability cost that is normalised by the
robot's own limits (slope limit and roughness saturation), embedded in A* as a
multiplicative edge cost with an admissible heuristic (\TA), followed by a
shortcutting step that provably never increases the integrated cost (\TAS).
Our main contribution is a careful, reproducible evaluation: \nTest{} test
scenarios over four procedurally generated terrain families and three map
scales, a held-out validation split used only to select the terrain weight,
planner-independent safety and energy metrics, and paired non-parametric
statistics. Against terrain-unaware \Astar{} and Theta*, \TAS{} reduces the
median traversability cost along the path by \costRelTAS\% relative to \Astar{}
and never enters terrain steeper than the slope limit, whereas \Astar{} and
Theta* do so in \exposedA\% and \exposedTheta\% of the scenarios; the median
path-length overhead is \lenRelTAS\% and the median change in estimated energy is
\energyRelTAS\%. Compared with a binary-traversability baseline (hard slope
constraint only), the continuous cost further lowers slope and roughness
exposure. We also report the trade-off over the terrain weight, a cost
ablation, robustness to elevation-map noise and computation time. Code and
data generators are released.
\end{abstract}

\begin{IEEEkeywords}
path planning, traversability, field robotics, unmanned ground vehicles, uneven terrain
\end{IEEEkeywords}

\section{Introduction}
Field robots in agriculture, construction and disaster response operate on
terrain where the shortest path is often not the best one: a straight line
over a hill or through a rough patch costs energy, increases the risk of
slipping or tipping, and may exceed the vehicle's climbing capability.
Terrain-aware planning addresses this by attaching a traversability cost to
the map and searching for paths that trade distance against terrain
difficulty~\cite{papadakis2013survey,guastella2021review}.

The idea itself is well established---planetary rovers have used
traversability maps for two decades~\cite{goldberg2002gestalt,ferguson2006fielddstar}---and
adding a terrain term to a graph-search cost is, algorithmically, a small
step. What is often less clear in practice is \emph{how much} such a cost
actually buys on a given robot, under which conditions, and at what price
in distance, energy and computation. Answering this requires evaluations
whose metrics are tied to the robot (e.g., distance driven on slopes beyond
its limit), whose baselines include both terrain-unaware and binary
traversability planners, whose parameters are not tuned on the test data,
and whose statistics respect the paired design of planner benchmarks.

This paper provides such an evaluation for a tracked UGV. Our contributions are:
\begin{enumerate}
\item A fully specified terrain-aware planner: a slope/roughness cost
normalised by the robot's limits rather than by per-map maxima, combined
with hard constraints, used as a multiplicative edge cost in A* with an
admissible heuristic (\TA), and a cost-aware shortcutting step (\TAS) that
never increases the integrated cost (Sec.~\ref{sec:method}).
\item A reproducible benchmark of \nTest{} test and \nVal{} validation
scenarios over four terrain families and three scales, with a pre-specified
rule for selecting the terrain weight on the validation split only
(Sec.~\ref{sec:setup}).
\item An evaluation with planner-independent safety and energy metrics and
paired, multiplicity-corrected statistics, including the distance--terrain
trade-off, an ablation of the cost terms, robustness to elevation noise and
computation time (Sec.~\ref{sec:results}).
\end{enumerate}
All numbers in this paper are generated from the released raw results by a
single script.

\section{Related Work}
\textbf{Grid search.} A*~\cite{hart1968astar} finds minimum-cost paths on
grids given an admissible heuristic; D* Lite~\cite{koenig2002dstarlite}
supports efficient replanning. Any-angle planners such as
Theta*~\cite{nash2007theta,nash2010lazy} and Field D*~\cite{ferguson2006fielddstar}
remove the grid's heading restriction; Field D* additionally interpolates
costs across cells and was deployed on Mars rovers. Sampling-based planners
such as RRT*~\cite{karaman2011rrtstar} scale to higher-dimensional state spaces
but are not the focus here.

\textbf{Traversability.} Surveys~\cite{papadakis2013survey,guastella2021review}
distinguish geometric approaches, which derive slope, roughness and step
height from elevation maps or point clouds, from appearance- and
learning-based ones. Elevation mapping with uncertainty~\cite{fankhauser2018elevation}
and SLAM systems such as RTAB-Map~\cite{labbe2019rtabmap} provide the maps on
which such features are computed. Geometric traversability costs have been
used for legged robots~\cite{wermelinger2016legged}, for planning directly on
point clouds~\cite{krusi2017pointclouds} and for energy-efficient paths on
uneven terrain~\cite{ganganath2015energy}.

\textbf{Positioning.} We do not propose a new search algorithm. We contribute
a carefully specified cost and planner for tracked UGVs and, mainly, an
evaluation protocol whose metrics, baselines, data split and statistics make
the benefit of terrain awareness measurable and reproducible.

\section{Method}\label{sec:method}
\subsection{Elevation map and terrain features}
The environment is a 2.5D elevation map $z(x,y)$ on a grid of resolution
$\Delta=0.5$\,m, as produced, for example, from an RTAB-Map point
cloud~\cite{labbe2019rtabmap}. For each cell we compute the slope
$\theta=\arctan\|\nabla z\|$ (central differences) and the roughness $\rho$,
the local standard deviation (3$\times$3 window) of the residual
$z - G_\sigma * z$ after removing smooth relief with a Gaussian filter
($\sigma=1.5$ cells). Geometric obstacles (rocks) are inflated by the robot's
footprint radius $r_f$.

\subsection{Robot-grounded cost and constraints}
Let $\thmax$ be the robot's admissible slope and $\rho_{\mathrm{sat}}$ a
roughness at which the roughness penalty saturates. The cost of a cell is
\begin{equation}
c = \alpha\,\min\!\Big(\frac{\theta}{\thmax},1\Big) + \beta\,\min\!\Big(\frac{\rho}{\rho_{\mathrm{sat}}},1\Big),\quad \alpha+\beta=1,
\label{eq:cost}
\end{equation}
so $c\in[0,1]$. Because the normalisers are robot properties, the same terrain
receives the same cost on every map; normalising by the per-map maximum, as
is sometimes done, makes costs incomparable across maps and lets a single
outlier rescale the whole map. Cells with $\theta>\thmax$ or inside an
inflated obstacle are \emph{lethal}. We use $\thmax=25^\circ$,
$\rho_{\mathrm{sat}}=0.10$\,m, $r_f=0.5$\,m and $(\alpha,\beta)=(0.7,0.3)$; these
values were fixed before any experiment and are ablated in Sec.~\ref{sec:ablation}.

\subsection{Terrain-Aware A* (\TA)}
The grid is 8-connected without corner cutting. The cost of moving between
neighbouring free cells $u,v$ at horizontal distance $d_{uv}$ is
\begin{equation}
k(u,v) = d_{uv}\Big(1 + w\,\frac{c_u + c_v}{2}\Big),
\label{eq:edge}
\end{equation}
with terrain weight $w\ge 0$. The path cost $g$ accumulates $k$; the
heuristic is the Euclidean distance to the goal.
\begin{proposition}
For any $w\ge 0$ the Euclidean heuristic is consistent for \eqref{eq:edge};
hence \TA{} returns a minimum-cost path.
\end{proposition}
\begin{proof}
Since $c\ge 0$, $k(u,v)\ge d_{uv}\ge h(u)-h(v)$ by the triangle inequality.
Consistency implies optimality of A*~\cite{hart1968astar}.
\end{proof}
Note that the terrain term must enter the \emph{accumulated} cost $g$: adding
it only to the priority $f=g+h+w\,c(n)$ changes the expansion order but not
the objective that the returned path optimises. With $w=0$, \TA{} reduces to A*
with hard constraints (\HC), i.e., binary traversability.

\subsection{Cost-aware shortcutting (\TAS)}
Grid paths contain heading artefacts. Standard string-pulling removes them
but can cut straight across costly terrain. \TAS{} shortcuts greedily: from
waypoint $p_i$ it extends a straight segment to the farthest $p_j$ such that
(i) the segment is collision-free with respect to lethal cells (a conservative
visibility test that also checks cells touched at boundaries) and (ii) its
integrated cost $\int(1+w\,c)\,ds$ does not exceed the cost of the original
sub-path $p_i\ldots p_j$ evaluated with the same quadrature.
\begin{proposition}
\TAS{} never increases the integrated path cost and never enters a lethal cell.
\end{proposition}
\begin{proof}
Each accepted shortcut replaces a sub-path by a segment of no greater cost
that passes the visibility test; the result is a concatenation of such
segments.
\end{proof}

\section{Experimental Setup}\label{sec:setup}
\subsection{Terrains and scenarios}
We generate four terrain families (Fig.~\ref{fig:examples}): \emph{hills}
(fractal relief, median slope $11^\circ$), \emph{ridge} (gentle relief with a
steep ridge across the start--goal line and one or two passes),
\emph{craters} (gentle relief with 6--11 crater-like depressions) and
\emph{rocky} (gentle relief, rough patches and scattered rocks). Relief is
produced by $1/f^\beta$ spectral synthesis and scaled to a target median slope,
so slope statistics are comparable across the three map sizes (50, 100 and
200\,m, i.e., up to 400$\times$400 cells). A centimetre-scale micro relief is
added everywhere. The start lies in the left and the goal in the right band
of the map. A scenario is kept only if a path exists under the hard
constraints, so every planner faces a solvable instance; \nRedraw{} of the
\nTest{} test scenarios needed a redraw. Each (family, size) cell contains
\nSeeds{} test seeds; a disjoint set of seeds (\nVal{} scenarios) forms the
validation split. All randomness is seeded and independent of the Python
hash seed.

\subsection{Planners}
\Astar{} (shortest path, obstacles only), Theta* (Lazy
Theta*~\cite{nash2010lazy}, any-angle shortest path, obstacles only), \HC{}
(shortest path with hard constraints), \TA{} and \TAS. All planners share the
same grid, neighbourhood, data structures and language (Python 3, single
thread), so their computation times are comparable to each other but not to
optimised C++ implementations.

\subsection{Terrain-weight selection}
The terrain weight is selected on the validation split with a rule fixed in
advance: the largest $w\in\{0,0.5,1,2,4,8,16,32\}$ whose median path-length
overhead of \TAS{} relative to $w=0$ is at most 10\%. This gives
$w=\wMain$ (median overhead \valOverheadMain\% on validation). The test split
is never used for tuning; the full sweep on the test split is reported for
transparency (Sec.~\ref{sec:sweep}).

\subsection{Metrics}
All paths are resampled every 0.1\,m and evaluated on the ground-truth map
with the same code: horizontal length; cumulative climb (sum of positive
elevation changes); maximum and 95th-percentile slope under the path;
\emph{steep length}, the distance driven on cells with $\theta>\thmax$;
mean cost $\bar c$ \eqref{eq:cost}; 95th-percentile roughness; and an energy
estimate $E = m g(\mu_r L + H^+)$ with $m=150$\,kg, rolling-resistance
coefficient $\mu_r=0.15$, length $L$ and climb $H^+$ (no regeneration). The
energy model is deliberately simple and is used only for relative
comparisons.

\subsection{Statistics}
Planners are compared on the same scenarios, so we use paired two-sided
Wilcoxon signed-rank tests~\cite{wilcoxon1945}, report the median paired
difference and the matched-pairs rank-biserial correlation $r$~\cite{kerby2014}
as effect size, and correct all \nTests{} tests with Holm's
method~\cite{holm1979}. Distributions are summarised by median and
interquartile range (IQR) because computation times and several metrics are
strongly skewed.

\section{Results}\label{sec:results}
% RESULTS-PLACEHOLDER

\end{document}
